%% Canonical pointwise (co)limits and componentwise map criteria.

\section{Presheaf categories and subobjects}

\begin{annexitem}{annex:p10.psh.1}{thm:p10-presheaf-limits}{Canonical pointwise limits and colimits}
\begin{proof}
\paragraph{The size of a hom-set.}
For fixed presheaves \(X,Y\), a natural transformation is a
member of
\[
                 \prod_{c\in\operatorname{Ob}\mathcal C}
                               \mathsf{Set}(X(c),Y(c))
\]
whose components satisfy the naturality equation for every
arrow of \(\mathcal C\).
The indexing object set, all function sets, and the resulting
product are small in the fixed universe.
The equations select a subset of that product.
Thus this actual natural-transformation set is small.
Here smallness of \(\mathcal C\), rather than only its local
smallness, is the hypothesis controlling the family of components.

\paragraph{Limits.}
Write \(D_j=D(j)\).
For each object \(c\), take the canonical set limit
\[
 L(c)=
 \left\{(x_j)_{j\in\operatorname{Ob}J}:
   x_j\in D_j(c),\
   D(a)_c(x_j)=x_k\text{ for every }a:j\to k\right\}.
\]
This is the compatible-tuples construction of
Theorem~\ref{thm:p10-limits-products-equalizers}.
For \(u:d\to c\), define
\[
                         L(u)((x_j)_j)=(D_j(u)(x_j))_j.
\]
The output tuple is compatible.  Indeed, naturality of
\(D(a):D_j\to D_k\) gives
\[
 D(a)_d D_j(u)(x_j)
        =D_k(u)D(a)_c(x_j)=D_k(u)(x_k).
\]
Identity and composition hold coordinatewise because they
hold for every \(D_j\).
Thus \(L\) is a presheaf.
Its coordinate projections \(\ell_j:L\to D_j\) are natural
by the definition of \(L(u)\), and form a cone.

For a cone \(\alpha_j:Z\to D_j\), define
\[
                       h_c(z)=(\alpha_{j,c}(z))_j.
\]
The cone equations make this tuple a member of \(L(c)\).
For \(u:d\to c\), naturality of each \(\alpha_j\) identifies
the components of \(L(u)h_c\) and \(h_d Z(u)\);
hence \(h\) is natural.
It has \(\ell_j h=\alpha_j\).
Any other mediator must have these same coordinates at
each object and element, and so equals \(h\).
This proves the limit universal property.

\paragraph{Colimits.}
For each \(c\), take the disjoint union
\[
                 B(c)=\coprod_{j\in\operatorname{Ob}J}D_j(c)
\]
and quotient it by the least equivalence relation identifying
\((j,x)\) with \((k,D(a)_c(x))\) for every \(a:j\to k\).
Call the quotient \(K(c)\).
For \(u:d\to c\), define its restriction on classes by
\[
                         K(u)([(j,x)])=[(j,D_j(u)(x))].
\]
For a generating pair, the two images are equivalent because
\[
                  D_k(u)D(a)_c(x)=D(a)_d D_j(u)(x).
\]
More formally, the relation of having the same image class
is an equivalence relation on \(B(c)\) containing each
generating pair.  It therefore contains the generated
equivalence relation.  The formula descends to the quotient
without choosing a representative for each class.
The identity and composition laws hold on every represented
element by the corresponding laws of \(D_j\), and hence
hold on the quotient.
Thus \(K\) is a presheaf.

The functions \(x\mapsto[(j,x)]\) give natural maps
\(k_j:D_j\to K\); the defining relations make them a cocone.
For any cocone \(\beta_j:D_j\to Z\), define
\[
                          h_c([(j,x)])=\beta_{j,c}(x).
\]
The cocone equations make this well-defined by the same
quotient argument.  Naturality follows from
\(\beta_{j,d}D_j(u)=Z(u)\beta_{j,c}\), evaluated on each
representative.
Every quotient element has a representative, so the
conditions \(h k_j=\beta_j\) also prove uniqueness.
These constructions establish the colimit universal property.

\paragraph{Empty cases and the particular constructions.}
When \(J\) is empty, the compatible tuple is the unique
empty tuple, and its set is a singleton at every object.
The empty disjoint union and its quotient are empty at
every object.  Their restrictions are the unique functions,
giving terminal and initial presheaves.
Discrete two-object diagrams give the stated product and
coproduct; the parallel-pair diagrams give the stated
equalizer subset and coequalizer quotient.
Equivalent choices of these set constructions give the
same universal properties through their canonical isomorphisms.

If \(\mathcal C\) is empty, there is one object-assignment
function, one arrow-assignment function, and one family of
transformation components, all empty.
There is therefore one presheaf and one natural arrow.
Every construction above is its unique object and arrow.
No pointwise limit, quotient, or inverse required an
independent family of choices.
\end{proof}
\end{annexitem}

\begin{annexitem}{annex:p10.psh.2}{thm:p10-presheaf-mono-epi}{Cancellation, images, and canonical subobject representatives}
\begin{proof}
\paragraph{Monomorphisms.}
If each \(m_c\) is injective and \(m\alpha=m\beta\) for
two natural maps into \(X\), their component equations
give \(\alpha_c=\beta_c\) by injectivity.
Thus \(\alpha=\beta\), proving that \(m\) is monic.

Conversely suppose \(m\) is monic and
\(m_c(x)=m_c(x')\) for \(x,x'\in X(c)\).
By Yoneda, these elements give natural maps
\(\alpha,\beta:\mathrm y c\to X\), with
\[
                    \alpha_d(u)=X(u)(x),\qquad
                    \beta_d(u)=X(u)(x')
                         \quad(u:d\to c).
\]
Naturality of \(m\) gives \(m\alpha=m\beta\) because their
values are \(Y(u)(m_c(x))\) and \(Y(u)(m_c(x'))\).
Monicity gives \(\alpha=\beta\).
Evaluate at \(1_c\) to get \(x=x'\), proving component
injectivity.

\paragraph{Epimorphisms.}
If each \(m_c\) is surjective and \(\alpha m=\beta m\),
then for any \(y\in Y(c)\) there exists \(x\in X(c)\)
with \(m_c(x)=y\).
The component equation gives
\(\alpha_c(y)=\alpha_c(m_c(x))=\beta_c(m_c(x))=\beta_c(y)\).
This proves \(\alpha=\beta\) without choosing a family
of such preimages.

Conversely let \(m\) be epic.
Take the pushout of \(X\xrightarrow{m}Y\) with itself,
with canonical arrows \(j_0,j_1:Y\to K\).
They satisfy \(j_0m=j_1m\), so epimorphy gives \(j_0=j_1\).
By Theorem~\ref{thm:p10-presheaf-limits}, this is a
pushout of sets at each component.
If some \(y\in Y(c)\) were outside the image of \(m_c\),
define \(a,b:Y(c)\to\{0,1\}\) by
\[
                  a(z)=0,\qquad
                  b(z)=
                    \begin{cases}1&z=y,\\0&z\ne y.\end{cases}
\]
Then \(a m_c=b m_c\).
The set pushout gives a function \(h:K(c)\to\{0,1\}\)
with \(h(j_{0,c}(z))=a(z)\) and
\(h(j_{1,c}(z))=b(z)\).
At \(z=y\) these are \(0\) and \(1\), contradicting
\(j_{0,c}=j_{1,c}\).
Thus every \(m_c\) is surjective.
The test function \(h\) is needed only for this set
component; it need not be a natural transformation.

\paragraph{Isomorphisms.}
An inverse natural transformation gives an inverse
function at every component.
Conversely suppose all \(m_c\) are bijective, and let
\(n_c=m_c^{-1}\) be their unique inverses.
For \(u:d\to c\), the naturality equation
\(Y(u)m_c=m_d X(u)\), composed with the two inverses,
gives
\[
                              X(u)n_c=n_dY(u).
\]
Thus the inverse functions form a natural transformation.
Their uniqueness defines the entire inverse family
without a choice of alternatives.
This proves the isomorphism criterion, and consequently
the assertion about maps that are both monic and epic.

\paragraph{Images and subobjects.}
For \(u:d\to c\), naturality gives
\[
             Y(u)(m_c(x))=m_d(X(u)(x))\in m_d[X(d)].
\]
Hence the image subsets are closed under restrictions
and form a subpresheaf \(I=\operatorname{Im}(m)\).
The component maps \(q_c:X(c)\to I(c)\) given by \(m_c\)
are surjective, and the inclusions \(i_c:I(c)\to Y(c)\)
are injective.  Their definitions and the displayed
equation make both families natural.
The proved criteria make \(q\) epic and \(i\) monic,
and \(m=iq\).
If \(m\) is monic, each \(q_c\) is also injective,
so \(q\) is an isomorphism.

Let \(m:X\to Y\) and \(n:Z\to Y\) be monic.
If their images agree, their just-constructed isomorphisms
onto that common image give a unique isomorphism
\(h:X\to Z\) with \(nh=m\).
Conversely, such an isomorphism is componentwise bijective
and makes the image subsets equal.
If the image of \(m\) is merely contained in that of \(n\),
compose the isomorphism \(X\to\operatorname{Im}(m)\),
the inclusion into \(\operatorname{Im}(n)\), and the
inverse isomorphism \(\operatorname{Im}(n)\to Z\).
This gives \(h\) with \(nh=m\).
Conversely any such factorization gives componentwise
image containment.  Its uniqueness follows from monicity
of \(n\).

Every subpresheaf inclusion is monic, so all families in
\(\operatorname{Sub}(Y)\) occur as representatives.
The set bound was established in
Definition~\ref{def:p10-subpresheaf}.
It concerns this set of canonical representatives, not
fixed-universe membership of every possible isomorphic
presentation of a subobject.
All the arguments remain valid for empty component sets
and for the empty base category.
\end{proof}
\end{annexitem}
